%==============================================================================
%  MATH 231 -- Linear Algebra I (Queens College, Fall 2026)
%  PLAN / NOT FINALIZED -- Topics that are NOT guaranteed to be covered.
%
%  Two groups:
%    (1) topics that remain in the progression document, tagged Recommended or
%        Enrichment, and are gathered here so the full set is visible at once;
%    (2) topics MIGRATED out of Unit 12, which now names them only under
%        "If time permits".
%==============================================================================
\documentclass[11pt]{article}

\usepackage[margin=1in]{geometry}
\usepackage{amsmath,amssymb}
\usepackage[dvipsnames]{xcolor}
\usepackage{enumitem}
\usepackage{booktabs}
\usepackage{array}
\usepackage{longtable}
\usepackage[hidelinks]{hyperref}
\usepackage{titlesec}

\titleformat{\section}{\Large\bfseries\sffamily\color{RoyalBlue}}{}{0pt}{}
\titlespacing*{\section}{0pt}{18pt}{8pt}

\title{\vspace{-1.2cm}\textbf{\sffamily MATH 231 -- Linear Algebra I}\\[4pt]
       {\large\sffamily Topics Not Guaranteed to Be Covered}\\[2pt]
       {\normalsize\sffamily Fall 2026 \textbullet\ working draft --- not finalized}}
\author{}
\date{}

\begin{document}
\maketitle
\vspace{-1.3cm}

\begin{center}
\fcolorbox{BrickRed!50}{BrickRed!7}{%
\begin{minipage}{0.93\textwidth}
\small\sffamily
\textbf{This document is not finalized.} Every topic listed here is one the
course \emph{may} reach, not one it promises to. Nothing here is part of the
syllabus and none of it should be studied in expectation of assessment unless it
has been officially announced. Where this document and an official course
document disagree, the official document is correct.
\end{minipage}}
\end{center}

\bigskip
\noindent The course progression is deliberately broader than one semester, so
that topics can be selected by interest and available time rather than fixed in
advance. This document collects, in one place, everything across the four arcs
that falls outside the guaranteed core --- so the boundary between
\emph{will cover} and \emph{might cover} is visible at a glance instead of
scattered across twelve units.

%==============================================================================
\section{Across the arcs}
%==============================================================================
These remain in the progression document, tagged \emph{Recommended} or
\emph{Enrichment}. They are gathered here because a tag buried in a unit is easy
to miss.

{\renewcommand{\arraystretch}{1.45}
\begin{longtable}{@{}p{4.4cm} c p{9.0cm}@{}}
\toprule
\textbf{Topic} & \textbf{Unit} & \textbf{What it would add} \\
\midrule
\endfirsthead
\toprule
\textbf{Topic} & \textbf{Unit} & \textbf{What it would add} \\
\midrule
\endhead
Quaternions & 3 &
  Non-commutative multiplication and 3-D rotation composition; the natural
  sequel to complex numbers as an algebra of rotations. \\
$LU$ with partial pivoting & 7 &
  Existence/uniqueness of the factorization once row interchanges are allowed;
  the practical form of Gaussian elimination. \\
Discrete Cosine Transform & 10 &
  A \emph{fixed} orthonormal basis for compression, contrasted against the
  data-adaptive SVD basis; the transform JPEG actually uses. \\
Graph Laplacians \& spectral clustering & 11 &
  $L = D - W$, and clustering a graph by the eigenvectors of its Laplacian. \\
Stochastic (Markov) matrices & 11 &
  Column-stochastic matrices and the dominant-eigenvector reading of a
  steady state. \\
Perron--Frobenius theorem & 11 &
  Why a nonnegative/stochastic matrix has a positive dominant eigenvector ---
  the result that makes PageRank well posed. \\
Eckart--Young & 11 &
  The truncated SVD is the best low-rank approximation; the theorem underneath
  every compression demo. \\
Randomized SVD & 11 &
  Sketching a large matrix to approximate its leading singular subspace. \\
Hutchinson trace estimation & 11 &
  Estimating $\operatorname{tr}(A)$ from matrix--vector products alone. \\
\bottomrule
\end{longtable}}

%==============================================================================
\section{Migrated out of Unit 12}
%==============================================================================
Unit~12 (\emph{Numerical Optimization \& Machine-Learning Capstone}) previously
carried the following as unit content. They have been moved here, and the unit
now names them only under \textbf{If time permits}. What Unit~12 still commits
to is gradient descent, Newton's method, and Gauss--Newton/Levenberg--Marquardt
for nonlinear least squares.

\begin{center}
\renewcommand{\arraystretch}{1.45}
\begin{tabular}{@{}p{4.4cm} p{9.6cm}@{}}
\toprule
\textbf{Topic} & \textbf{What it would add} \\
\midrule
Quasi-Newton \textbf{BFGS} &
  Building curvature information from successive gradients instead of forming
  and solving with the Hessian directly. \\
\textbf{Support vector machines} &
  Constrained optimization under the Vapnik objective; the maximum-margin
  classifier, and kernel methods reusing the similarity kernels from Unit~2.
  Unit~2 retains a geometric preview of the SVM margin. \\
Stochastic gradient descent (\textbf{SGD}) &
  Descent on noisy estimates of the gradient; the method that makes
  large-scale fitting tractable. Also underlies the masked-loss / ALS
  factorization behind the collaborative-filtering demo. \\
Model-selection methodology &
  Train/validation/test splits and cross-validation as estimates of
  generalization error. \\
\bottomrule
\end{tabular}
\end{center}

\bigskip
\noindent\textbf{Dependency worth noting.} Candidate project \#7
(\emph{Mini-ML: OLS + SVM + KNN with validation}, in
\texttt{MATH231\_projects\_and\_demos}) draws on both the support vector machine
and model-selection methodology. If those topics are not reached, that project
is not assignable in its current form.

\end{document}
